% ============================================================================
%  Semantic Uncertainty Propagation
%  A self-contained theory of individuation, meaning, and comprehension-free
%  translation in finite weighted graphs.
%
%  Self-contained: every object is a finite weighted graph or a construction
%  on one; only the standard external literature is cited. Interpretation is
%  confined to remarks on which no theorem depends.
% ============================================================================
\documentclass[twocolumn,11pt,a4paper]{article}

\usepackage[utf8]{inputenc}
\usepackage[T1]{fontenc}
\usepackage{lmodern}
\usepackage[a4paper,margin=2cm]{geometry}
\usepackage{amsmath,amssymb,amsthm,mathtools}
\usepackage{bm}
\usepackage{mathrsfs}
\usepackage{booktabs}
\usepackage{array}
\usepackage{enumitem}
\usepackage{microtype}
\usepackage{xcolor}
\usepackage[round,sort&compress]{natbib}
\usepackage[colorlinks=true,linkcolor=blue!45!black,
            citecolor=blue!45!black,urlcolor=blue!45!black]{hyperref}
\usepackage{cleveref}

% ---- theorem environments ----
\newtheoremstyle{sup}{6pt}{6pt}{\itshape}{}{\bfseries}{.}{.5em}{}
\theoremstyle{sup}
\newtheorem{theorem}{Theorem}[section]
\newtheorem{lemma}[theorem]{Lemma}
\newtheorem{proposition}[theorem]{Proposition}
\newtheorem{corollary}[theorem]{Corollary}
\theoremstyle{definition}
\newtheorem{definition}[theorem]{Definition}
\newtheorem{axiom}{Premise}
\newtheorem{construction}[theorem]{Construction}
\newtheorem{example}[theorem]{Example}
\theoremstyle{remark}
\newtheorem{remark}[theorem]{Remark}
\newtheorem{notation}[theorem]{Notation}

\crefname{axiom}{Premise}{Premises}
\Crefname{axiom}{Premise}{Premises}

% ---- operators and shorthand ----
\DeclareMathOperator{\Aut}{Aut}
\DeclareMathOperator*{\argmin}{arg\,min}
\DeclareMathOperator*{\argmax}{arg\,max}
\newcommand{\R}{\mathbb{R}}
\newcommand{\N}{\mathbb{N}}
\newcommand{\Z}{\mathbb{Z}}
\newcommand{\Gph}{\Gamma}
\newcommand{\V}{V}
\newcommand{\E}{E}
\newcommand{\wt}{w}
\newcommand{\med}{\mathsf{m}}
\newcommand{\flo}{\beta}
\newcommand{\cut}{\partial}
\newcommand{\co}{\complement}
\newcommand{\str}{\sigma}
\newcommand{\res}{\rho}
\newcommand{\rec}{M}
\newcommand{\al}{a}
\newcommand{\Om}{\Omega}
\newcommand{\dem}{D}
\newcommand{\colA}{A}
\newcommand{\colB}{B}
\newcommand{\Resp}{\mathsf{R}}
\newcommand{\rlx}{\Phi}
\newcommand{\abs}[1]{\lvert#1\rvert}
\newcommand{\norm}[1]{\lVert#1\rVert}
\newcommand{\eps}{\varepsilon}
\newcommand{\dd}{\,\mathrm{d}}

\title{\textbf{Semantic Uncertainty Propagation}\\[3pt]
\large A Self-Contained Theory of Individuation, Meaning,\\
and Comprehension-Free Translation in Finite Weighted Graphs}

\author{%
Kundai Farai Sachikonye\\[2pt]
\small Technical University of Munich\\
\small Chair of Experimental Bioinformatics\\
\small AIMe Registry for Artificial Intelligence\\
\small \texttt{kundai.sachikonye@wzw.tum.de}}

\date{\today}

\begin{document}
\maketitle

% ============================================================================
\begin{abstract}
\noindent
We develop, from a single derived fact, a complete theory of meaning and its
transfer. The fact is a strictly positive \emph{floor} $\flo>0$ on the cost of
telling one thing from the rest of a whole that cannot be finished: identifying
a part is comparing it against everything it is not, and against a
non-completable whole that comparison never closes, so an irreducible residue
remains. Every object below is a finite weighted graph or a construction on
one. From the floor we prove, as graph theorems: that a part is fixed only by
its negation (\emph{individuation by negation} admits no selector without
regress); that the \emph{identity} of a part is not a label but the invariant
minimum cut that survives every relabelling, a region of weight $\ge\flo$ and
never a point---so identity is the \emph{least sufficient unknowability} of a
thing, never the thing itself; that a \emph{name} is the content-free handle on
this cut---a stash of unknowability---and the \emph{self}-name is necessarily late
and approximate, the namer (having the largest complement) being the most expensive
thing to individuate and forbidden by the diagonal of self-verification from naming
itself directly, so that language begins at the individuated periphery and converges
on the namer it never reaches; that an item must be \emph{terminated} (carry an
attained resting cut) before it can be told apart at all; that the \emph{content}
of a unit (its label) and its \emph{meaning} (its resting cut, and what that
cut propagates) come apart, and that meaning cannot be extracted---reading it
off as a value would require the reader to occupy the part's whole
negation-structure, i.e.\ to \emph{be} it. We therefore do not try to know
meaning; we detect when two sides \emph{share} it. We replace every imported
uncertainty calculus---probability, fuzzy membership, Bayesian posterior---by a
single parameter-free quantity, \emph{distance to the floor}, and define a
\emph{bidirectional relaxation} of two columns whose fixed point, \emph{semantic
quiescence}, we prove to \emph{constitute} identity of meaning: any difference
of meaning is a propagable demand and quiescence is the absence of all demands,
so the identity is certified \emph{blind}, at positive residual of form, never
at zero; non-halting is deep untranslatability, whose honest output is to
decline. We then prove that meaning is receiver-relative with no privileged
value (distinct decoders register distinct cells, each correct), that a
\emph{conversation} is propagation of uncertainty in which a response is forced
exactly because no point-meaning can be extracted and is steered by the demand
rather than by comprehension, and we construct the \emph{four-column
translator}: two central translation columns and two external
\emph{response} columns that externalise the propagation, so that translation
becomes the test of whether two utterances start the same conversation---an
audit of the route, not the endpoints, which a meaning-blind mechanism passes
and a false friend fails. A master equivalence binds every structure to the one
fact $\flo>0$. The development is offered as the structural common form beneath
the indeterminacy of translation, radical interpretation, conversational
implicature, and meaning-as-use; we prove everything of the abstract graph
alone, and confine interpretation to remarks.
\end{abstract}

\noindent\textbf{Keywords:} finite weighted graph; minimum cut; individuation
by negation; resolution floor; meaning as region; content versus meaning;
semantic quiescence; bidirectional relaxation; distance to the floor;
comprehension-free translation; propagation of uncertainty; conversation.

\tableofcontents

% ============================================================================
\part{Foundations: the floor}
% ============================================================================

% ============================================================================
\section{Introduction}
\label{sec:intro}
% ============================================================================

\subsection{The question}

What is the meaning of an utterance, and what is it to render it faithfully in
another tongue? Two traditions answer differently. One locates meaning in
\emph{reference} and \emph{sense}---a sentence means by picking out conditions
in the world \citep{frege1892}---and translation is the preservation of those
conditions. Another locates meaning in \emph{use}---a word means by its role in
a practice \citep{wittgenstein1953}---and translation is the matching of roles.
Each has been pressed to a limit: that translation is \emph{indeterminate}, with
no fact of the matter fixing one rendering over an incompatible other
\citep{quine1960}; that interpretation is \emph{radical}, recoverable only
holistically from behaviour \citep{davidson1973}; that what is communicated
outruns what is said, by \emph{implicature} read off the assumption that a
speaker is cooperating \citep{grice1975,sperberwilson1986}; and that meaning
``ain't in the head'' \citep{putnam1975}. We do not adjudicate among these. We
exhibit a single mathematical structure of which each is a projection, and we
derive it---not assume it---from one fact about finitely many parts told apart
within a whole they cannot exhaust.

\subsection{The thesis}

Three claims organise the paper, each proved below.

\emph{First}, telling a part from the rest is never free. In a finite weighted
graph of parts individuated within a non-completable whole, every cut separating
a part from the rest carries weight bounded below by a strictly positive floor
$\flo>0$ (\Cref{sec:floor}). The floor is derived: a complete identification of
a part would have to bring the whole to bear, which a non-completable whole
forbids, so a residue remains.

\emph{Second}, the identity of a part is that residue---the invariant minimum
cut, a \emph{region} of weight $\ge\flo$, never a point (\Cref{sec:identity}). It
is the \emph{least sufficient unknowability} of a thing: the cheapest commitment
of \emph{what it is not} that suffices to tell it apart. It is never the thing
itself; to reach the thing itself (residue $0$) one would have to be it.
Meaning, accordingly, cannot be extracted (\Cref{sec:meaning}).

\emph{Third}, because meaning cannot be read, identity of meaning is detected,
not known: two sides mean the same iff no difference can \emph{propagate}
between them. We make propagation a graph dynamics, define its fixed point
(\emph{quiescence}), and prove that quiescence \emph{constitutes} identity of
meaning---certified blind, at positive residual of form
(\Cref{sec:quiescence}). The rest of the paper turns this into a translator
that never comprehends and a model of conversation in which a response is forced
by the residue and steered by the demand (\Cref{sec:conversation,sec:four}).

\subsection{Method}

The development is austere on purpose: every object is a finite weighted graph
or a construction on one, with no metric, no measure, and no probability except
as a derived normalisation. The reason is to leave nothing undefined. ``The
identity of a thing'' is a specific graph invariant; ``meaning'' is a specific
cut; ``uncertainty'' is a specific distance; ``a proper translation'' is a
specific fixed point. A reader who grants only finite weighted graphs can check
every claim, and minimum cuts are computable exactly by maximum
flow~\citep{fordfulkerson1956,menger1927}. Where a tempting claim could not be
proved at this standard it has been removed rather than weakened, and every
interpretive reading is confined to a remark on which no theorem depends.

% ============================================================================
\section{Contact graphs}
\label{sec:graph}
% ============================================================================

\begin{definition}[Finite weighted graph; cut]
\label{def:graph}
A \emph{finite weighted graph} is a triple $\Gph=(\V,\E,\wt)$ with $\V$ a finite
non-empty vertex set, $\E\subseteq\binom{\V}{2}$ undirected edges, and
$\wt\colon\E\to\R_{>0}$. We take $\Gph$ connected unless stated. For
$\varnothing\neq S\subsetneq\V$ the \emph{cut} is
$\cut(S)=\{\{u,v\}\in\E:u\in S,\ v\notin S\}$ with weight
$\wt(\cut(S))=\sum_{e\in\cut(S)}\wt(e)$.
\end{definition}

\begin{definition}[Contact graph; the medium]
\label{def:contact}
A \emph{contact graph} is a finite weighted graph with a distinguished vertex
$\med\in\V$, the \emph{medium}, adjacent to every other vertex. The vertices
$v\neq\med$ are the \emph{items}; an edge between two items is a \emph{contact},
its weight the cost of the separation distinguishing them.
\end{definition}

\begin{definition}[Separation cost]
\label{def:sepcost}
The \emph{separation cost} of an item $v$ is the minimum weight of a cut placing
$v$ on one side and the medium on the other,
\[
\str(v)\ =\ \min_{S\ni v,\ \med\notin S}\ \wt(\cut(S)),
\]
attained on finite $\V$; the minimising cut is the \emph{resting cut} of $v$.
This is the $v$--$\med$ minimum cut, computable by max-flow
\citep{fordfulkerson1956,menger1927,diestel2017}.
\end{definition}

\begin{remark}[The medium is the rest, not a background]
\label{rem:medium}
The medium is not a spatial container; it is the single vertex standing for
``everything else,'' against which each item is told apart. Its adjacency to
every item encodes that every item is, at minimum, individuated against the
whole. No theorem uses any reading of $\med$ beyond \Cref{def:contact}.
\end{remark}

% ============================================================================
\section{The floor, derived}
\label{sec:floor}
% ============================================================================

We show the central constant is not assumed but forced by two facts: an item is
a proper part, and the whole is non-completable.

\begin{definition}[Expanding family; non-completable whole]
\label{def:noncompletable}
An \emph{expanding family} is a chain of contact graphs
$\Gph_0\subseteq\Gph_1\subseteq\cdots$, each obtained from the last by adding
items and contacts while retaining the medium and all prior structure. The whole
is \emph{non-completable} if the family has no terminal stage: for every finite
$\Gph_k$ there is $\Gph_{k+1}\supsetneq\Gph_k$. This is an order condition (no
last stage), not a cardinality assumption.
\end{definition}

\begin{definition}[Identification residual]
\label{def:residual}
To \emph{identify} an item $v$ at stage $\Gph_k$ is to determine it by its
resting cut $\cut(S^\ast(v))$ in $\Gph_k$. The identification is \emph{complete}
if every item of the whole has been brought to bear in the cut---i.e.\ the cut
is taken in a graph containing the whole. The \emph{identification residual}
$\res(v)\ge0$ is the cut content a complete identification would remove that a
finite-stage identification cannot; $\res(v)=0$ iff identification is complete.
\end{definition}

\begin{theorem}[Floor from infinitude]
\label{thm:floor-derived}
Let the whole be non-completable and let $v$ be a proper part ($v\neq\med$, $v$
not exhausting $\V$). Then the identification of $v$ cannot be completed at any
finite stage, and $\res(v)>0$.
\end{theorem}

\begin{proof}
Suppose $\res(v)=0$: identification is complete at some finite $\Gph_k$. By
\Cref{def:residual} completeness means every item of the whole has entered the
cut of $v$ against $\med$, i.e.\ the cut is computed in a graph containing the
whole, so $\Gph_k$ contains the whole. This contradicts non-completability
(\Cref{def:noncompletable}), which guarantees a further stage
$\Gph_{k+1}\supsetneq\Gph_k$ with an item not in $\Gph_k$, hence not yet brought
to bear. No finite stage completes the identification; $\res(v)>0$.
\end{proof}

\begin{theorem}[Positive floor]
\label{thm:floor}
There is a constant $\flo>0$ with $\wt(e)\ge\flo$ for every contact and
$\str(v)\ge\flo$ for every item: no two items are separated at zero cost, and no
item is individuated for free.
\end{theorem}

\begin{proof}
Set $\flo:=\inf_v\res(v)$. Each $\res(v)>0$ (\Cref{thm:floor-derived}); the
residual is the irreducible boundary content of $v$'s individuation, carried by
the edges of $v$'s resting cut, so each such edge has weight $\ge\flo$. Every
edge $\{u,v\}$ lies in $\cut(\{u\})$, so every weight is $\ge\flo$. Connectivity
forces $\cut(S^\ast(v))\neq\varnothing$, so
$\str(v)=\wt(\cut(S^\ast(v)))\ge\flo>0$.
\end{proof}

\begin{corollary}[No sharp cut]
\label{cor:no-sharp}
There is no cut of weight $0$ separating a connected contact graph into a
nonempty item set and its nonempty complement. Every individuation deposits
weight $\ge\flo$.
\end{corollary}

\begin{remark}[The floor is the trace of the whole on its parts]
\label{rem:trace}
\Cref{thm:floor-derived,thm:floor} make $\flo$ derived: it is the residue a
non-completable whole forces on every finite part. A completable whole would
permit $\res(v)=0$, hence sharp cuts and point-identities; the impossibility of
finishing the comparison is exactly what makes the floor positive. We return to
this as the master equivalence (\Cref{sec:master}); no later theorem assumes
$\flo$, only its positivity.
\end{remark}

% ============================================================================
\begin{figure*}[t]
\centering
\includegraphics[width=\textwidth]{figures/panel_1.png}
\caption{\textbf{The floor and granular meaning} (computed on finite weighted
graphs; exact minimum cuts). \textbf{(A)}~Separation cost $\sigma(v)$ for every
item of ${\sim}220$ random contact graphs versus graph size; every point lies on
or above the floor $\flo$ (red), independent of size---no sharp cut
(\Cref{thm:floor}). \textbf{(B)}~Distribution of $\sigma(v)/\flo$, supported
entirely at and above $1$. \textbf{(C)}~Alignment score $\al(x,y)$ versus the
floor $\flo/\Om$; every point lies above the diagonal, the bound
$\al\ge\flo/\Om$ (\Cref{lem:align-floor}). \textbf{(D)}~The realised floor
$\min_v\sigma(v)$ over graph size and $\flo$: strictly positive throughout,
rising with $\flo$.}
\label{fig:panel1}
\end{figure*}

\section{Individuation by negation}
\label{sec:negation}
% ============================================================================

\begin{definition}[Selector; no privileged vertex]
\label{def:selector}
A \emph{selector} for an item set $U\subseteq\V$ is a rule fixing $U$ by naming
a distinguished vertex of $U$ in advance. A determination is \emph{selector-free}
if it names no vertex of $U$. The graph has \emph{no privileged vertex} if the
data presented to a determining rule are invariant under the weighted
automorphism group $\Aut(\Gph)$, so no vertex is singled out by the data alone.
\end{definition}

\begin{theorem}[Individuation by negation]
\label{thm:negation}
In a contact graph with no privileged vertex, any selector-free determination of
an item set $U$ fixes $U$ as the complement of its negation set,
$U=\V\setminus\co U$ with $\co U:=\V\setminus U$ specified first; conversely
every admissible $\co U$ determines a unique $U$; and no selector-based rule
fixes $U$ without invoking a further selector.
\end{theorem}

\begin{proof}
\emph{Necessity.} To fix $U$ is to decide membership of each vertex. A positive
rule asserting membership of a distinguished vertex must name it; with
$\Aut(\Gph)$-invariant data no vertex is named, so the rule must supply
one---a selector---which must in turn be fixed, either given in advance
(excluded) or by a prior rule facing the same dichotomy. The regress halts only
at a rule naming no vertex of $U$: the withholding of $U$ from the whole, i.e.\
specifying $\co U$ and taking $U$ as the remainder, which names only ``the
rest,'' supplied by $\V$. \emph{Sufficiency and uniqueness.} Complementation
$U\mapsto\V\setminus U$ is an involution on subsets of the finite set $\V$, so
$\co U$ determines $U$ uniquely and conversely.
\end{proof}

\begin{corollary}[An item is fixed by its non-neighbours]
\label{cor:nonneighbours}
A single item $v$ is determined selector-free by the partition of the rest into
its neighbours $N(v)$ (what it contacts) and non-neighbours $\V\setminus N[v]$
(what it neither is nor contacts), with no intrinsic mark on $v$. Its identity
is carried by this negation structure, not by its label.
\end{corollary}

\begin{remark}[Why negation closes the regress]
\label{rem:regress}
A positive selector is a rule ``choose a vertex'' that must itself be chosen;
with no privileged vertex the choice has no ground and regresses
(\Cref{def:selector}). Complementation is the unique operation fixing a part
using only the whole and the part's own absence---data already present once the
graph is given. Negation, not selection, is the primitive.
\end{remark}

% ============================================================================
\section{Identity is the invariant cut: the least sufficient unknowability}
\label{sec:identity}
% ============================================================================

We identify what about an item is invariant across all re-encodings of the
graph, and show it is never a label and never a point.

\begin{definition}[Weighted isomorphism; relabelling]
\label{def:iso}
A \emph{weighted isomorphism} $\phi\colon\Gph\to\Gph'$ is a vertex bijection
preserving the medium, edges, and weights. A \emph{relabelling} is an
automorphism. A \emph{graph invariant} is a function equal on isomorphic graphs.
\end{definition}

\begin{theorem}[The minimum cut is the conserved invariant]
\label{thm:invariant}
The separation cost $\str(v)$ and the resting cut of $v$ are weighted-graph
invariants: if $\phi\colon\Gph\to\Gph'$ is a weighted isomorphism then
$\str_{\Gph'}(\phi v)=\str_\Gph(v)$. The vertex label of $v$ is not invariant: it
is permuted by automorphisms.
\end{theorem}

\begin{proof}
$\phi$ carries each $v$--$\med$ cut to a $\phi v$--$\med'$ cut of equal weight
(weights preserved), and the correspondence on cuts has inverse $\phi^{-1}$, so
the minima agree. Labels are permuted by any nontrivial automorphism, hence not
invariant.
\end{proof}

\begin{theorem}[Identity is a region, never a point]
\label{thm:region}
The identity of an item---the invariant datum fixing it---is its resting cut, a
set of contacts of total weight $\str(v)\ge\flo>0$. It is realised by no single
vertex: in general the minimising side $S^\ast(v)$ is not the singleton
$\{v\}$, and the identity cannot be reduced to a zero-content (point) datum.
\end{theorem}

\begin{proof}
By \Cref{thm:invariant} the invariant fixing $v$ is the minimum-cut value and
its cut set, not the variable label. By \Cref{thm:floor} the cut has weight
$\ge\flo>0$, never a point. Non-locality: take two dense item-clusters joined by
a single $\flo$-weight contact; the minimum cut against $\med$ may separate a
whole cluster (a side of size $>1$) at weight below $\min_v\wt(\cut(\{v\}))$, so
the minimiser is not a singleton and the identity is borne by a region of
contacts.
\end{proof}

\begin{definition}[Least sufficient unknowability]
\label{def:lsu}
The \emph{least sufficient unknowability} of an item $v$ is its separation cost
$\str(v)$: the minimum commitment of \emph{what $v$ is not} that suffices to
tell $v$ from the rest. It is the cheapest sufficient individuation, attained at
the resting cut, of weight $\ge\flo$.
\end{definition}

\begin{theorem}[Identity is the least sufficient unknowability, never the thing]
\label{thm:lsu}
For every item $v$: \textup{(i)} $\str(v)$ is the least cost at which $v$ is
sufficiently individuated (a smaller cut does not separate $v$ from $\med$);
\textup{(ii)} $\str(v)\ge\flo>0$, so the identity is never the thing itself (a
zero-cost, residue-free datum); \textup{(iii)} attaining residue $0$---``being
$v$ itself''---would require completing the comparison of $v$ against the whole,
forbidden by \Cref{thm:floor-derived}.
\end{theorem}

\begin{proof}
(i) By \Cref{def:sepcost} $\str(v)$ is the minimum $v$--$\med$ cut; any cheaper
cut fails to separate $v$ from $\med$, so does not individuate $v$. (ii) is
\Cref{thm:floor}. (iii) residue $0$ is $\res(v)=0$, which
\Cref{thm:floor-derived} shows requires the whole to be brought to bear, i.e.\
the comparison completed, impossible against a non-completable whole.
\end{proof}

\begin{remark}[A name denotes a cut]
\label{rem:name}
A name---a label such as ``this thing''---carries no invariant; it is the
relabellable surface. What is invariant is the resting cut, the region of
contrasts by which the thing is told from the rest, of weight $\ge\flo$. A name
denotes that cut at the current stage of an expanding family, converging to an
intrinsic value it never reaches as a point (\Cref{thm:region,thm:lsu}). This is
the precise sense in which the identity of a thing is what it is \emph{not},
sufficiently and at least cost---never what it is.
\end{remark}

% ============================================================================
\begin{figure*}[t]
\centering
\includegraphics[width=\textwidth]{figures/panel_2.png}
\caption{\textbf{Individuation by negation and the invariant region.}
\textbf{(A)}~The minimum-cut minimiser side $|S^\ast(v)|$ versus cluster size $k$
for the two-cluster construction: it equals $k>1$, so identity is a region, never
a point (\Cref{thm:region}). \textbf{(B)}~$\sigma(v)$ after a random relabelling
versus before: on the diagonal---identity is a relabelling invariant
(\Cref{thm:invariant}). \textbf{(C)}~$|U|+|\complement U|$ versus $|V|$: on the
diagonal, a part and its negation partition the whole exactly
(\Cref{thm:negation}). \textbf{(D)}~The same in three coordinates
$(|U|,|\complement U|,|V|)$, lying on the plane $|U|+|\complement U|=|V|$.}
\label{fig:panel2}
\end{figure*}

\part{Termination, meaning, and propagation}
% ============================================================================

% ============================================================================
\section{Termination: only completed items individuate}
\label{sec:termination}
% ============================================================================

Before a part can be told apart there must be a part to tell. We isolate the
condition---\emph{termination}---and show it is prior to everything that
follows.

\begin{definition}[Terminated and open units]
\label{def:terminated}
A \emph{unit} is a connected item set $U$ with the induced contact structure. It
is \emph{terminated} if its resting cut against the medium is attained, giving a
definite separation cost $\str(U)\ge\flo$; it is \emph{open} if no such cut is
attained (the minimising cut does not exist as a committed boundary).
\end{definition}

\begin{theorem}[Only terminated units can be told apart or propagated]
\label{thm:terminate}
A unit has a defined alignment to another unit (\Cref{def:alignment} below),
hence a defined demand to propagate (\Cref{def:demand}), only if terminated. An
open unit has no defined cut against the medium, hence no defined alignment and
no computable demand.
\end{theorem}

\begin{proof}
Alignment and demand are defined through a minimum cut against the medium
(\Cref{def:alignment,def:demand}), which requires an attained resting
cut---termination. An open unit's separation cost is undefined; no alignment or
demand is computable on it.
\end{proof}

\begin{construction}[Completion]
\label{con:completion}
A unit is terminated by committing a completion boundary: a bounded
(``completed'') form, drawn as a cut, gives the resting cut the rest of the
theory requires. Some completion must be applied for a unit to enter the
development; which one is a modelling choice. \textit{(No theorem depends on the
choice.)}
\end{construction}

\begin{remark}[Termination is the first step, and it is not free]
\label{rem:first-step}
\Cref{thm:terminate} marks termination as the mandatory first operation: an
un-terminated ``distinction'' commits no boundary, hence individuates nothing
and propagates nothing---there is no part there to tell apart. Committing the
completion boundary deposits residue $\ge\flo$ (\Cref{cor:no-sharp}); the act of
making a thing a definite thing already pays the floor. In linguistic terms this
is casting an utterance into a completed, well-formed unit before it can be
rendered; nothing downstream applies to an item that has not first been brought
to a boundary.
\end{remark}

% ============================================================================
\section{Content, meaning, and the impossibility of extraction}
\label{sec:meaning}
% ============================================================================

\begin{definition}[Content versus meaning]
\label{def:content}
The \emph{content} of a unit is its label---its literal marks and their
dictionary correspondents. The \emph{meaning} of a unit $U$ is its resting cut
(\Cref{thm:region}): the region of contrasts individuating $U$ against the rest.
Content is a name on the cell; meaning is the cell's individuation against the
whole. They coincide only when the literal form is itself the operative
individuation, which is not the general case (\Cref{ex:greeting}).
\end{definition}

\begin{example}[Content and meaning come apart]
\label{ex:greeting}
``How are you?'' has content \emph{health-inquiry} and, in greeting usage,
meaning \emph{greeting}: its resting cut individuates it against other
greetings, not against other health-inquiries. A rendering that preserves
content yields the health-inquiry correspondent and misses the greeting;
the two land on different target cells \citep{jakobson1959,nida1964}.
\end{example}

\begin{theorem}[Cognate substitution cannot reach meaning]
\label{thm:no-cognate}
Let two contact graphs $\Gph,\Gph'$ (two languages) carve the same space of
senses with different weights. A content map---one sending each unit to its
correspondent label---does not in general preserve meaning: there are units
whose correspondent has a different resting cut from the unit's meaning. Two
obstructions guarantee it: \textup{(i)} \emph{different carvings} (one graph
resolves a cell where the other resolves several, so a source meaning need not
be any single target cell); \textup{(ii)} \emph{content $\neq$ meaning} (even
within one carving, the correspondent preserves content, which is not the
meaning).
\end{theorem}

\begin{proof}
(i) Let $\Gph$ resolve one cell where $\Gph'$ resolves two; the content map
sends the source cell to one correspondent, whose resting cut individuates it
against the \emph{other} target cell---a contrast absent from the source cell's
cut---so the cuts differ. (ii) In \Cref{ex:greeting} the correspondent preserves
the health-inquiry content; its resting cut individuates against
health-inquiries while the unit's meaning individuates against greetings. In
both cases the correspondent's meaning differs from the source meaning.
\end{proof}

\begin{theorem}[Meaning is unextractable]
\label{thm:unextractable}
There is no operation that reads the meaning of a unit off as a value---i.e.\ that
resolves its resting cut to a residue-free point. Such a reading would require
$\res(U)=0$, which by \Cref{thm:floor-derived} requires the whole to be brought
to bear, i.e.\ the reader to occupy $U$'s entire negation-structure---to
\emph{be} $U$, which is no longer a comparison. Hence no agent, in either
language, reads meaning as a value.
\end{theorem}

\begin{proof}
By \Cref{thm:region,thm:lsu} the meaning is the resting cut of weight $\ge\flo$,
and reducing it to a point is $\res(U)=0$, excluded by
\Cref{thm:floor-derived}. Concretely, to fix $U$ exactly is to decide it against
the totality of what it is not; the agent reading it is itself within that
totality, so a complete reading would require it to be both the part read and
the whole it is read against---to occupy both sides of the constituting
partition at once. The gap that keeps them apart is the residue $\ge\flo$.
\end{proof}

\begin{remark}[The precise reading is a precise reading of content]
\label{rem:precision}
Optimising a rendering for precision optimises fidelity to the source
\emph{content}---the literal marks and their correspondents. By
\Cref{def:content} content is exactly what the meaning is not. A precise
rendering therefore reproduces the non-meaning faithfully: an accurate
\emph{contact} that is not the unit's significance. Precision is fidelity to the
wrong object; the right object is the meaning, which is what the unit
\emph{propagates} (\Cref{sec:propagation}), and which \Cref{thm:unextractable}
forbids reading. This motivates the criterion of \Cref{sec:quiescence}; no
theorem depends on the reading.
\end{remark}

\begin{remark}[The standing form beneath indeterminacy]
\label{rem:quine}
\Cref{thm:unextractable} is the structural root of the indeterminacy of
translation \citep{quine1960} and the holism of radical interpretation
\citep{davidson1973}: there is no fact reading meaning off a word as a value,
because the value would be a residue-free point the floor forbids. What we add
is that indeterminacy is not a defect to be lamented but the condition under
which there is anything to say (\Cref{sec:conversation}); and that, though
meaning cannot be read, identity of meaning can be \emph{detected}, which is all
a translator needs (\Cref{sec:quiescence}).
\end{remark}

% ============================================================================
\section{Names: stashes of unknowability, and the late self-name}
\label{sec:names}

One unit deserves separate treatment: the \emph{name}. We show a name is the
content-free handle on an item's resting cut---a stash of unknowability---and
that the \emph{self}-name, the name of the namer, is necessarily late and
approximate, read directly off individuation by negation together with a
self-contained diagonal.

\subsection{A name is the cut without the content}

\begin{definition}[Name]
\label{def:name}
A \emph{name} of an item $v$ is a token $n_v$ that fixes $v$'s cell---its resting
cut $\mathrm{cell}(v):=\cut(S^\ast(v))$---by reference, adding no contrast of its
own to that cut. A name individuates by pointing; it carries no descriptive
content.
\end{definition}

\begin{theorem}[A name is a stash of unknowability, never the thing]
\label{thm:name-stash}
For an item $v$ with name $n_v$: \textup{(i)} $n_v$ individuates $v$ (fixes a
unique cell) while carrying no descriptive content (it adds nothing to $v$'s
resting cut); \textup{(ii)} what $n_v$ retrieves is $v$'s cell---the region of
contrasts of weight $\ge\flo$ by which $v$ is told from the rest---never $v$
itself, the residue-free invariant, which is unextractable
(\Cref{thm:unextractable}); \textup{(iii)} hence a name is the irreducible
minimum individuation of $v$---a floor-handle---packaging $v$'s unknowability
($\ge\flo$ of \emph{not}-$v$) into a single referenceable token.
\end{theorem}

\begin{proof}
(i) By \Cref{def:name} $n_v$ fixes the cell (individuates) without adding any
contrast (no content). (ii) The cell is the resting cut, a region of weight
$\ge\flo$ (\Cref{thm:region}), the retrievable datum; $v$ itself is the
residue-$0$ point, forbidden and unextractable
(\Cref{thm:lsu,thm:unextractable}). (iii) The cell is the least sufficient
unknowability of $v$ (\Cref{def:lsu}); $n_v$ is its handle.
\end{proof}

\begin{remark}[The two routes; a name versus a description]
\label{rem:two-routes}
There are two selector-free ways to fix $v$: \emph{unpack} the resting cut
explicitly (describe $v$ against everything it is not), or \emph{reference} it by
a name. A proper name---``John''---is the limit: it makes a \emph{specific} $v$
(maximal individuating force) while adding no description (zero content), a token
denoting a cut without connoting it. It does not convey what $v$ is; it stashes
\emph{what $v$ is not}, sufficiently, in one word. And what it returns is the
retrievable cell, not $v$: what one retrieves is what one could figure out, not
the thing. No theorem depends on the reading.
\end{remark}

\subsection{Static and dynamic irreducibles; a name is not reached by negation}

\begin{definition}[Static and dynamic irreducibles]
\label{def:static-dynamic}
A \emph{static irreducible} is a resting cut individuated context-free---the same
cell in every context (the catalogue cell of a fixed item, an element). A
\emph{dynamic irreducible}, or \emph{name}, is a stash whose referenced cell is
fixed only relative to a running context: the cell it individuates against shifts
as the context shifts.
\end{definition}

\begin{theorem}[A name is not reached by a negation-sequence]
\label{thm:non-derivable}
No name is the terminus of a negation-sequence. \textup{(i)} A name-token is a
handle (\Cref{def:name}), not a cell, so it is not in the range of the cut
operation: no walk of committing cuts outputs the name; the name is \emph{given},
not computed. \textup{(ii)} A name is \emph{non-compositional}: the resting cut of
a compound name need not equal any composition of its constituents' resting cuts,
so no walk of negations over the constituents reaches the compound's cell.
\end{theorem}

\begin{proof}
(i) By \Cref{def:name} a name references a cell externally; the cut operation
(\Cref{def:sepcost}) produces a cell---an edge set---never a token, so the token
is not its output. (ii) Let $U=\{u,v\}$ be a compound whose items are densely
joined to each other (weight $W$) but thinly to the rest (weight $\flo$ each).
Its resting cut is the minimum $U$--$\med$ cut, of weight $2\flo$ and cut set
$\{u\med,v\med\}$; the would-be composition of the constituents' singleton cuts
is $\{uv,u\med\}\cup\{uv,v\med\}$, of weight $W+2\flo$ and a different edge set.
For $W>0$ these differ, so the compound's cell is not the composition of the
parts' cells, and no negation over the parts reaches it.
\end{proof}

\begin{remark}[The riverbed]
\label{rem:riverbed}
A description is the negation-sequence---``not this, not that,'' terminating at a
cell---reachable by such steps. A name short-circuits it: ``riverbed'' is one
element whose cell is \emph{not} ``a bed where a river sleeps''; ablating either
word destroys the individuation, and the literal sub-senses do not compose to it
(\Cref{thm:non-derivable}\,(ii)). One cannot \emph{derive} a name; it is handed
over. This is why a name is a stash and not a label on a known thing: a label on
a known thing would be computable, and a name is not. No theorem depends on the
reading.
\end{remark}

\begin{proposition}[Contextual interchangeability]
\label{prop:interchange}
In a fixed context, distinct handles referencing the same cell---a name, a
description, a circumlocution---have the same alignment to every target, hence the
same cross-demand, hence are quiescent with one another (\Cref{thm:quiescence});
the particular handle is immaterial, and only the referenced cell enters any
invariant.
\end{proposition}

\begin{proof}
Alignment is a minimum-cut weight, a function of the referenced cell and not of
the handle (\Cref{def:alignment}); two handles on the same cell give equal
alignment to every target, so all cross-demands agree and they are quiescent
(\Cref{thm:quiescence}). The handle, a label, is permuted by relabelling and is
not an invariant (\Cref{thm:invariant}).
\end{proof}

\begin{theorem}[A dynamic stash is resolved only by propagation]
\label{thm:stash-propagation}
The cell a name stashes cannot be recovered by decomposing the name---it is not
reached by negation (\Cref{thm:non-derivable})---and is context-relative
(\Cref{def:static-dynamic}, \Cref{thm:receiver}). The only handle-independent,
context-sensitive quantity is the propagation---the cross-demand, equivalently the
response the name invites. Hence identity of two stashes is identity of their
propagation, tested by the four-column criterion (\Cref{thm:four}), never by
decomposition.
\end{theorem}

\begin{proof}
By \Cref{thm:non-derivable} the stashed cell is not the terminus of a derivation
from the token; by \Cref{thm:receiver} it is decoder/context-relative; so no
handle-internal or context-free quantity fixes it. The cross-demand
(\Cref{def:demand}) is handle-independent (\Cref{prop:interchange}) and
context-sensitive, and its vanishing is identity of meaning (\Cref{thm:quiescence});
so identity of stash is identity of propagation, which \Cref{thm:four} tests by
the responses.
\end{proof}

\begin{remark}[Element versus name; rest versus propagation]
\label{rem:element-name}
A static irreducible is a cell \emph{at rest}---an element of a catalogue, the
same in every context (``lithium'' always means lithium). A name is a cell
\emph{under propagation}---fixed only relative to a context, and resolved only
through how it propagates. ``Peter,'' ``a medical student at Harvard,'' and ``one
who studies human anatomy at an institution of higher learning in Massachusetts''
are, in a context, one dynamic stash under three handles; which handle is used is
immaterial (\Cref{prop:interchange}), and that they stash the same thing is
established by their starting the same conversation (\Cref{thm:stash-propagation}),
not by any decomposition. No theorem depends on the reading.
\end{remark}

\subsection{The self-name is the diagonal}

\begin{definition}[Namer; self-naming]
\label{def:namer}
A \emph{namer} is an agent, realised as a part of the graph, that given an item
$X$ and a candidate constituting totality $T$ returns $\mathsf{V}(X,T)\in\{0,1\}$,
reading $1$ iff $X$ is the part $T$ constitutes, i.e.\ iff $X=\V\setminus T$. To
\emph{name the namer} is to evaluate $\mathsf{V}$ on the part its own verdicts
constitute---its own negation.
\end{definition}

\begin{theorem}[No direct self-name: the diagonal]
\label{thm:self-diagonal}
Let $D:=\{X:\mathsf{V}(X,\co X)=0\}$. If $\mathsf{V}$ is total (returns a value on
every part and its negation) and $D$ is itself a part on which
$\mathsf{V}(\cdot,\co\,\cdot)$ is defined, then $\mathsf{V}(D,\co D)$ has no
consistent value: $\mathsf{V}(D,\co D)=1$ and $\mathsf{V}(D,\co D)=0$ each entail
the other. Hence no namer names itself directly; the self-name cannot be
assigned at the source.
\end{theorem}

\begin{proof}
If $\mathsf{V}(D,\co D)=1$ then $D$ is declared constituted by $\co D$; but the
membership condition of $D$ is $\mathsf{V}(X,\co X)=0$, so $D$'s satisfying its
own condition forces $\mathsf{V}(D,\co D)=0$. If instead $\mathsf{V}(D,\co D)=0$
then $D$ meets its membership condition, hence is among the parts $D$ collects,
which is the assertion $\mathsf{V}(D,\co D)=1$. Both values are self-refuting;
none is consistent. This is the diagonal of \citet{cantor1891,russell1903}, the
undecidability shape of \citet{godel1931,tarski1936}, and the fixed-point form of
\citet{lawvere1969}.
\end{proof}

\begin{remark}[On the closure hypothesis]
\label{rem:closure}
The theorem is exactly as strong as its hypothesis that the behaviour-defined $D$
is again a part on which $\mathsf{V}$ acts---genuine \emph{self}-application. An
agent that names while being itself a part of the graph makes this natural: its
verdicts leave a trace in the same whole they range over. We do not claim to
discharge it in general; we mark it, and note that the unextractability of
\Cref{thm:unextractable} (the floor argument) holds without it. What rides on the
hypothesis is the sharper claim that the self-name is forbidden \emph{directly},
hence reached only terminally.
\end{remark}

\subsection{The self-name is late, hence an approximation}

\begin{theorem}[The self-name is late and approximate]
\label{thm:self-late}
A name is fixed by negation against the already-named (\Cref{thm:negation}). The
namer's negation is everything it has named---the largest complement---so its
separation cost is maximal and its individuation the most expensive. Therefore:
\textup{(i)} the self-name cannot be first---it requires a non-empty set of prior
names to be individuated against; \textup{(ii)} the direct self-name is forbidden
(\Cref{thm:self-diagonal}), so the self-name is reached only \emph{terminally},
after a periphery of other names creates the comparison; \textup{(iii)} it never
reaches the invariant (\Cref{thm:unextractable}), so it is an
\emph{approximation}---a floor-handle on the namer, not the namer.
\end{theorem}

\begin{proof}
By \Cref{thm:negation} an item is fixed only as the complement of its negation;
the namer's negation is the set of everything it has named, which must therefore
be non-empty and prior, giving (i). By \Cref{thm:self-diagonal} the namer cannot
take itself as its own input directly, so no early, direct self-name is
assignable; the only available determination is the indirect accumulation by
negation against the already-named, which is terminal once that periphery exists,
giving (ii). By \Cref{thm:unextractable,thm:name-stash} the handle never reduces
the cut to residue $0$, giving (iii).
\end{proof}

\begin{corollary}[Names approximate an unknowability; they are not labels on
known things]
\label{cor:names-approx}
Were a name a label on a \emph{known} thing, the best-known thing---the namer
itself, the source of all knowing---would be named first. It is named last. The
order is inverted from the label-on-a-known-thing prediction; hence names do not
radiate outward from the best-known but accumulate inward from the individuated
periphery, reaching toward an unknowability they never close. The self-name is
the terminal approximation of the hardest case.
\end{corollary}

\begin{proof}
A label on a known thing would be assignable first to the most-known; the namer
is the most-known yet named last (\Cref{thm:self-late}\,(i,ii)); the only
consistent reading is that the name approximates an unextractable invariant
(\Cref{thm:self-late}\,(iii)) rather than carrying it.
\end{proof}

\begin{remark}[Language begins at the periphery]
\label{rem:periphery}
Interpretively, \Cref{cor:names-approx} is the founding asymmetry of language:
the first word was a name for something that is \emph{not} the self, because one
can name the self only against an already-named everything-else, and the
everything-else must come first. The thing doing all the naming is the last to
receive a name, and the name it receives---``mind,'' ``brain,'' a vapour-handle
on the namer---is not it but a provisional floor-handle reaching toward the
invariant it can never extract. The self-name is the linguistic instance of the
diagonal: self-reference is the terminal operation, arrived at through everything
else, and even then only approximately. No theorem depends on the reading.
\end{remark}

% ============================================================================
\section{The residue is the propagator; the record is time}
\label{sec:propagation}
% ============================================================================

We make cuts dynamic: a process is a sequence of committed cuts, the residue of
each is what enables the next, and the count of cuts is the system's time.

\begin{definition}[Committed cut; process; record]
\label{def:process}
A \emph{cut event} commits a contact $e\in\E$ as a realised separation. A
\emph{process} is a finite sequence $P=(e_1,\dots,e_k)$ of distinct committing
acts. The \emph{committed record} is $\rec(P)=k$. The \emph{residue} of $e_i$ is
$\res_i:=\wt(e_i)\ge\flo>0$.
\end{definition}

\begin{theorem}[Residue as causal propagator]
\label{thm:propagator}
In a process the residue $\res_i>0$ of cut $e_i$ is necessary and sufficient for
$e_{i+1}$ to be a distinct committing event: $e_{i+1}$ is a new cut iff the state
after $e_i$ differs from the state before, which holds iff $\res_i>0$.
\end{theorem}

\begin{proof}
$e_{i+1}$ is distinct iff it acts on a state changed by $e_i$. The only change
$e_i$ makes is depositing its residue $\res_i$. If $\res_i>0$ the post-state
carries a committed contact absent before, so $e_{i+1}$ acts on a strictly
different configuration: distinct. If $\res_i=0$ (excluded by \Cref{thm:floor},
argued for necessity) pre- and post-states coincide and $e_{i+1}$ has no changed
state to act on. Thus $\res_i>0$ is necessary and sufficient.
\end{proof}

\begin{corollary}[A process is a residue chain]
\label{cor:chain}
A process is a chain in which each cut's residue causes the next:
$\res_1\to e_2,\ \res_2\to e_3,\dots$; remove any $\res_i$ and the chain halts
at step $i$.
\end{corollary}

\begin{theorem}[Monotone non-return]
\label{thm:monotone}
The committed record is strictly increasing along any process; no process
decreases it. A process revisiting a configuration does not return to a prior
state: the configuration recurs but the record has strictly grown, and the two
states are distinguished by the record.
\end{theorem}

\begin{proof}
Each cut event increments $\rec$ by one (\Cref{def:process}). To decrease $\rec$
a committed contact would have to be un-committed; un-committing is itself a
further committing act (a new cut redrawing the boundary), incrementing $\rec$.
Hence $\rec$ is monotone. Two states with equal configuration but records
$r<r'$ are separated by the partition ``record $\le r$'' vs ``$>r$,'' hence
distinct: recurrence of configuration is not return of state.
\end{proof}

\begin{remark}[The record is time]
\label{rem:time}
\Cref{thm:monotone} attaches a monotone scalar to every process---the committed
count---never decreasing, with equal configurations at different counts being
different states. ``Before'' and ``after'' are the order on the record;
the order is intrinsic to the sequence of cuts, requiring no external parameter.
\end{remark}

\begin{theorem}[Three readings of a residue]
\label{thm:three}
A residue $\res_i$ admits exactly three readings, by how its link to the next
cut is resolved: \textup{(i)} \emph{time}---as the increment of the record
(\Cref{thm:monotone}); \textup{(ii)} \emph{information}---when tracked as the
specific cause of a specific next cut, a resolved link; \textup{(iii)}
\emph{entropy}---when residues are present but their links are untracked, their
aggregate is the Shannon content $H=-\sum_i p_i\log p_i$ of the unresolved-link
distribution \citep{shannon1948,cover2006}. Tracking moves a process from
entropy toward information; losing track moves it back.
\end{theorem}

\begin{proof}
(i) is \Cref{thm:monotone}. (ii) a tracked link is a determinate map
``$\res_i\mapsto e_{i+1}$''; specifying it for each $i$ specifies the chain
exactly. (iii) with links unresolved the chain is known only through the
distribution over which residue caused which next cut, of information content
$H$ \citep{shannon1948}. Resolving links converts unknown distribution (entropy)
into determinate functions (information).
\end{proof}

% ============================================================================
\part{The uncertainty calculus and the relaxation}
% ============================================================================

% ============================================================================
\section{Distance to the floor: the parameter-free uncertainty}
\label{sec:distance}
% ============================================================================

We replace every imported uncertainty calculus by one derived quantity: how far
a state is from the floor.

\begin{definition}[Alignment; alignment score]
\label{def:alignment}
Fix a contact graph and a target unit $y$. The \emph{alignment} of a unit $x$ to
$y$ is the minimum-cut weight separating them,
$\str(x,y)=\min_{S\ni x,\ y\notin S}\wt(\cut(S))\in[\flo,\Om]$, where
$\Om:=\wt(\E)$ is the total weight. The \emph{alignment score} is
$\al(x,y):=\str(x,y)/\Om\in(0,1]$, small when $x,y$ are cheaply separated
(nearly the same cut), large when expensively separated.
\end{definition}

\begin{lemma}[The floor bounds alignment below]
\label{lem:align-floor}
For distinct units $x\neq y$, $\al(x,y)\ge\flo/\Om>0$: no two distinct units are
perfectly aligned; $\al=0$ would require $\str=0$, forbidden by \Cref{thm:floor}.
\end{lemma}

\begin{proof}
Immediate from \Cref{thm:floor}: any $x$--$y$ cut has weight $\ge\flo$.
\end{proof}

\begin{theorem}[Distance to the floor replaces the borrowed calculus]
\label{thm:replace}
The quantity ordinarily carried by a probability, a fuzzy membership, or a
Bayesian posterior---a number in $[0,1]$ recording how sure or how much---is
carried throughout by the alignment $\al(\cdot,\cdot)$ and its distance to the
floor $\flo/\Om$. This quantity is \textup{(i)} \emph{derived} (a minimum-cut
weight of the graph, not a fitted parameter), \textup{(ii)} \emph{bounded below}
by $\flo/\Om>0$ (no certainty to a point, \Cref{lem:align-floor}), and
\textup{(iii)} \emph{parameter-free} (no prior is set, no membership tuned).
A resolution is to within the floor, never to a point; a residual $\res\ge\flo$
always remains.
\end{theorem}

\begin{proof}
(i) $\al=\str/\Om$ with $\str$ an exact minimum cut (\Cref{def:alignment}). (ii)
\Cref{lem:align-floor}. (iii) the definition fixes $\al$ from the weights alone;
no external prior or membership function enters. The least achievable distance
is the floor itself, so resolution stops at $\flo/\Om>0$, never at $0$.
\end{proof}

\begin{remark}[Why a derived floor, not a prior]
\label{rem:bayes}
The standard calculi---subjective probability \citep{definetti1937},
its logic \citep{cox1946,jaynes2003}, fuzzy membership \citep{zadeh1965}---each
introduce a quantity by stipulation: a prior, a membership function, a
degree. Here the corresponding quantity is forced: it is the trace the
non-completable whole leaves on each part (\Cref{rem:trace}), so no prior is
set and no membership is tuned. Where an uncertainty would be \emph{posited},
the floor is \emph{derived}. No theorem depends on the comparison.
\end{remark}

% ============================================================================
\section{Cross-demand and bidirectional relaxation}
\label{sec:relaxation}
% ============================================================================

\begin{definition}[Two columns; anchors and residue]
\label{def:columns}
A \emph{column} is a terminated unit together with its current target rendering,
ablation-ranked into \emph{anchors}---the items whose removal most lowers the
unit's separation cost, held fixed---and \emph{residue items}---the rest, which
the relaxation may move. The \emph{source} is column $\colA$; a naive
correspondent rendering (each item to its label correspondent) is the
\emph{seed} column $\colB$.
\end{definition}

\begin{construction}[Noun anchoring]
\label{con:noun}
After termination (\Cref{con:completion}), take the \emph{nouns} of a unit as the
default anchor set. Each noun is a name---a committed cut, a context-relative
stashed cell (\Cref{thm:name-stash,def:static-dynamic})---so the material it
stashes is \emph{sufficiently excluded} (excluded to within the floor, relative to
the running context) and need not be propagated into; the
relaxation moves only the residue, the relational and predicative structure
\emph{among} the nouns. A noun is \emph{re-opened} (un-anchored) precisely when
its cross-demand is positive---when its stash does not transfer (a carving
mismatch, \Cref{thm:no-cognate}, or a false friend)---anchors being stable but
reclassifiable (\Cref{thm:relax-mono}).
\end{construction}

\begin{remark}[The nouns are the catalogue; the relations are the propagation]
\label{rem:noun-catalogue}
\Cref{con:noun} gives the relaxation a grammatical first pass: the nouns demarcate
a \emph{boundary of material not to consider}---already stashed, sufficiently
excluded---and the relaxation works only on what propagates, the structure among
them. ``The bank approved the loan'' and ``the loan approved the bank'' share
their noun-stashes but differ in relation, hence in meaning: the meaning lives in
the residue, not the stashes. A sentence is then its noun-stashes---\emph{dynamic}
irreducibles (\Cref{def:static-dynamic}), fixed only relative to the running
context---together with its propagation (the relations among them); the static
catalogue, the context-free cells, is the limiting element-like case
(\Cref{rem:element-name}). Exclusion is
\emph{sufficient}, never perfect: each noun keeps its own form-residual $\ge\flo$,
checked at quiescence. No theorem depends on the reading.
\end{remark}

\begin{definition}[Reverse propagation; cross-demand]
\label{def:demand}
A \emph{reverse propagation} of column $X$ into the other language is a walk in
the shared medium (the single uncarved whole both languages individuate, joined
by contacts wherever a cell of one is individuable against a cell of the other)
from each residue cell of $X$ to a cell of the other language, holding $X$'s
anchors fixed; each step commits residue $\ge\flo$ and the cause of the next step
is the residue deposited by the last. The \emph{cross-demand}
$\dem_{\colA\to\colB}$ is the change required to make $\colB$ consistent with
$\colA$'s reverse-image, measured as the summed above-floor alignment gap
\[
\dem_{\colA\to\colB}=\!\!\sum_{i\in\mathrm{res}(\colB)}\!\!
\bigl(\al(\colB_i,\rlx_{\colA\to\colB}(\colA)_i)-\tfrac{\flo}{\Om}\bigr)\ge0,
\]
and $\dem_{\colB\to\colA}$ symmetrically. A column is \emph{satisfied from the
other} when that direction's cross-demand is zero.
\end{definition}

\begin{definition}[Relaxation; quiescence; proper translation]
\label{def:relax}
The \emph{relaxation} alternates $\rlx_{\colA\to\colB}$ (update $\colB$'s residue
to satisfy $\colA$, then re-ablate) and $\rlx_{\colB\to\colA}$, and iterates. The
pair is \emph{quiescent}---a \emph{fixed point}---when both cross-demands vanish,
$\dem_{\colA\to\colB}=\dem_{\colB\to\colA}=0$. A translation is \emph{proper}
iff quiescent.
\end{definition}

\begin{theorem}[Monotone, floor-bounded relaxation]
\label{thm:relax-mono}
Along the relaxation the total committed residue is non-decreasing (each update
commits cuts $\ge\flo$ and none is withdrawn, by \Cref{thm:monotone}) and is
bounded below by $\flo$ per maintained contact; after the anchors stabilise
(finitely many updates, since each column has finitely many items) the
cross-demands are non-increasing.
\end{theorem}

\begin{proof}
Each update is a reverse propagation committing cuts $\ge\flo$; by monotone
non-return (\Cref{thm:monotone}) the committed total never decreases. Anchors,
being the finitely many maximal-ablation items, are re-classified at most until
their ranking stabilises, after finitely many updates. Thereafter each update
only moves residue toward satisfying the other side's demand, which by
\Cref{def:demand} reduces that direction's cross-demand or leaves it unchanged.
\end{proof}

\begin{theorem}[Dichotomy: quiescence or non-halting]
\label{thm:dichotomy}
The relaxation either reaches quiescence in finitely many steps, or it does not
halt. After anchor stabilisation the cross-demands are non-increasing and bounded
below by $0$, hence convergent; if the limit is $0$ over a finite state space the
infimum is attained at a finite step (quiescence); if the limit is a positive
constant the demands never vanish and the relaxation does not halt.
\end{theorem}

\begin{proof}
By \Cref{thm:relax-mono} the demands are non-increasing and $\ge0$, so
convergent. The state space (finitely many residue assignments over finitely
many cells) is finite, so a $0$ limit is attained at a finite step. A positive
limit is never reached, and each update still requires a change on the other
side. These are the only cases. The relaxation is a contraction toward the fixed
point or it has none, in the sense of \citep{banach1922,tarski1955}.
\end{proof}

% ============================================================================
\section{Quiescence constitutes identity of meaning}
\label{sec:quiescence}
% ============================================================================

\begin{theorem}[Quiescence is identity of meaning, certified blind]
\label{thm:quiescence}
If $(\colA,\colB)$ is quiescent then $\colA$ and $\colB$ have the same meaning,
and conversely a difference of meaning makes them non-quiescent; hence
\[
\text{quiescence}\ \Longleftrightarrow\ \text{identity of meaning},
\]
and the equivalence holds without any party knowing the meaning.
\end{theorem}

\begin{proof}
\emph{Quiescent $\Rightarrow$ identical.} Suppose quiescent yet differing in
meaning. A meaning-difference is a resting-cut difference: a contrast
individuating $\colA$ that does not individuate $\colB$, or conversely.
Reverse-translate $\colA$ into $\colB$ (\Cref{def:demand}); where $\colA$'s
meaning carries a contrast $\colB$ lacks, the reverse-image places a cell
$\colB$'s residue does not match, so
$\al(\colB_i,\rlx_{\colA\to\colB}(\colA)_i)>\flo/\Om$ there, whence
$\dem_{\colA\to\colB}>0$---contradicting quiescence. \emph{Difference
$\Rightarrow$ non-quiescent.} The same construction yields a positive
cross-demand. \emph{Blind.} The argument refers only to cross-demands, which are
alignment gaps---minimum-cut weights---computed without resolving any cell to a
value (\Cref{thm:unextractable}); the meaning is never read. Identity is
certified by the \emph{absence of a propagable difference}, not by inspection.
\end{proof}

\begin{theorem}[Form survives: quiescence at positive residual]
\label{thm:form}
A proper translation is quiescent at residual $\res\ge\flo>0$, never at
$\res=0$: the two sides have identical meaning and irreducibly different form,
$\res$ being the form-residue that does not cross-propagate at quiescence.
\end{theorem}

\begin{proof}
Every maintained contact carries weight $\ge\flo$ (\Cref{thm:floor}); the two
columns are distinct carvings (\Cref{thm:no-cognate}), so even at identity of
meaning each retains its own form-cells, individuated at cost $\ge\flo$ that the
other cannot absorb without becoming the same language. Thus the residual is
$\ge\flo>0$, and it is form, not meaning (cross-demands are zero). $\res=0$ would
be a zero-cost identity of expression, forbidden by \Cref{cor:no-sharp}.
\end{proof}

\begin{theorem}[Two untranslatabilities; the honest decline]
\label{thm:two-kinds}
A unit fails to translate in exactly two ways, distinguished by the relaxation:
\textup{(i)} \emph{shallow}---the relaxation halts but no single target cell is
at floor distance, resolved at quiescence by a composite (gloss) or a borrowing
(a difference of carving); \textup{(ii)} \emph{deep}---the relaxation does not
halt, so no fixed point exists and the carvings admit no mutual semantic
equilibrium. On a deeply untranslatable unit the correct output is to
\emph{decline}---to report that no proper translation exists---rather than emit a
non-quiescent (meaning-divergent) rendering.
\end{theorem}

\begin{proof}
(i) If quiescent, the meanings are identical (\Cref{thm:quiescence}); identity is
realised by a single floor-distance cell, or by the minimal composite whose joint
cell reaches floor distance, or by a borrowing when none does (decided by
computing single- and composite-cell alignments at the fixed point). (ii) If
non-halting, the cross-demands are bounded below by a positive constant
(\Cref{thm:dichotomy}); a nonzero cross-demand is a meaning-difference
(\Cref{thm:quiescence}), so a difference persists under every update---no
assignment makes the meanings identical. Any emitted rendering would be
non-quiescent, i.e.\ meaning-divergent; declining is the only output consistent
with properness.
\end{proof}

\begin{remark}[Why this is the criterion, not precision]
\label{rem:criterion}
A precision-seeker halts when the target matches the source content; the
relaxation halts when no meaning-difference can propagate (\Cref{rem:precision}).
``How are you?'' relaxes to the target greeting, not the health-inquiry
correspondent: reverse-translating the greeting demands no health-inquiry, so the
pair is quiescent at the greeting. The precision-seeker halts early, at content;
the relaxation halts at meaning. No theorem depends on the example.
\end{remark}

\begin{figure*}[t]
\centering
\includegraphics[width=\textwidth]{figures/panel_3.png}
\caption{\textbf{Propagation, the monotone record, and the relaxation.}
\textbf{(A)}~The committed record $M$ along a walk (blue), strictly increasing;
red marks revisited vertices---recurrence of position is not return of state
(\Cref{thm:monotone}). \textbf{(B)}~Cross-demand $D$ along the relaxation: one
instance reaches $0$ (quiescence, teal), one plateaus at the floor (non-halt
$\to$ decline, orange) (\Cref{thm:dichotomy,thm:two-kinds}). \textbf{(C)}~Residual
above floor for several seeds, each descending to $0$. \textbf{(D)}~The demand
surface $D(\text{step},D_0)$: monotone descent to the floor.}
\label{fig:panel3}
\end{figure*}

% ============================================================================
\part{Conversation and the four-column translator}
% ============================================================================

% ============================================================================
\section{Receiver-relativity: no privileged meaning}
\label{sec:receiver}
% ============================================================================

\begin{definition}[Receiver; registered cell]
\label{def:receiver}
A \emph{receiver} is a contact graph (a decoder of percepts into individuated
cells). The \emph{cell a unit registers} in a receiver is the cell of that
receiver's graph the unit lands in---its individuation against \emph{that}
receiver's rest.
\end{definition}

\begin{theorem}[Distinct receivers register distinct cells, each correct]
\label{thm:receiver}
The same unit registers different cells in receivers with different decoder
graphs, and each registration is a correct individuation at that receiver's
floor. No registration is privileged as ``the'' meaning; there is no
point-meaning of which the registrations are approximations.
\end{theorem}

\begin{proof}
By \Cref{def:receiver} the registered cell is the unit's resting cut in the
receiver's graph; different graphs give different cuts (a sparse graph
individuates the unit against few contrasts---a thin cell; a dense graph against
many---a rich cell). Each is the minimum-cut individuation in its own graph, so
each is a correct (floor-attaining) individuation there. Were there a privileged
point-meaning, a registration differing from it would be incorrect; but each
registration is the correct individuation in its own graph, and by
\Cref{thm:unextractable} no point-meaning exists to be approximated. Hence none
is privileged.
\end{proof}

\begin{remark}[The Harvard student]
\label{rem:harvard}
``A medical student at Harvard'' registers, in a small child's sparse graph, the
thin cell \emph{a student}; in a recruiter's dense graph it registers a rich
cell whose contrasts propagate---\emph{not} a dropout, \emph{not} below the top
percentile, \emph{not} idle. Both register a meaning; neither is the meaning. The
propagated contrasts are not in the content (``not necessary'') yet are present
(the residue $\ge\flo$ that propagates). This is conversational implicature
\citep{grice1975,sperberwilson1986} and meaning-externalism \citep{putnam1975}
in one structure: what is communicated is the receiver's individuation of the
unit against its own rest, not a value in the words. No theorem depends on the
reading.
\end{remark}

% ============================================================================
\section{Conversation as propagation of uncertainty}
\label{sec:conversation}
% ============================================================================

\begin{definition}[Conversation]
\label{def:conversation}
A \emph{conversation} between two receivers is a process (\Cref{def:process}) in
which each terminates an utterance into a cell, and the residue of each
utterance is the cause of the next. There is no ledger of cells to satisfy in
order and no pre-fixed terminal cell; the participants are uncertain and do not
hold in advance the cell they will reach.
\end{definition}

\begin{theorem}[A response is forced by the absence of a point-meaning]
\label{thm:response}
In a conversation, because no utterance's meaning can be extracted as a value
(\Cref{thm:unextractable}), the exchange cannot terminate at a known cell; a
response is required, and is exactly the non-vanishing cross-demand---the residue
not yet relaxed to quiescence (\Cref{def:relax}). Were a point-meaning available,
no response would be required (the value would already be held); its
unavailability is what forces the next utterance.
\end{theorem}

\begin{proof}
By \Cref{thm:unextractable} the utterance's meaning is not a value any party
holds. By \Cref{def:relax} two terminated cells are at rest only when both
cross-demands vanish; until then $\dem>0$, and reducing a positive cross-demand
requires a move on the other side---a response (\Cref{def:demand}). If instead a
point-meaning were held, the relaxation would have nothing to propagate and the
pair would be at rest with no response demanded; this is the excluded $\res=0$
case (\Cref{thm:form}). Hence a response is forced exactly by the positivity of
the residual.
\end{proof}

\begin{theorem}[Steering is by the demand, not the meaning]
\label{thm:steer}
The next utterance is determined by the propagation of uncertainty---the standing
cross-demand, i.e.\ what further individuation is required to reduce
it---and not by the meaning, which need not be read. Comprehension is optional:
the next move is a function of the residual gap alone.
\end{theorem}

\begin{proof}
By \Cref{def:relax} the only quantity governing the update is the cross-demand;
by \Cref{thm:unextractable} the meaning is not available to govern it, and by
\Cref{thm:quiescence} identity (and so the stopping condition) is itself read off
the demand, not the meaning. Hence the selection of the next utterance is a
function of the cross-demand---what is still required---and is independent of any
comprehended value. The demand is the selector of the next step
(\Cref{cor:nonneighbours}: what is required is read from the negation structure,
not from a label).
\end{proof}

\begin{remark}[Uncertainty is the condition for there being anything to say]
\label{rem:condition}
\Cref{thm:response} inverts the usual picture: uncertainty is not noise to be
removed on the way to communication but the condition of communication. A unit
with no residual ($\res=0$) propagates nothing, demands no response, and is the
forbidden floor-zero object (\Cref{cor:no-sharp}); an utterance that expects no
response is no utterance. The same floor that forbids perfect alignment
(\Cref{lem:align-floor}) is what makes there be anything to say. This is the
conversational reading of bounded rationality \citep{simon1955}: a finite
participant proceeds by propagating uncertainty, not by holding values it cannot
have.
\end{remark}

% ============================================================================
\section{The four-column translator}
\label{sec:four}
% ============================================================================

Meaning is unextractable, but its \emph{propagation} is observable: we cannot
read the meaning, but we can generate what a unit would say next---its
response---and responses are comparable across languages. We make this an extra
relaxation layer.

\begin{construction}[Four columns]
\label{con:four}
Given a source unit, form four columns:
\begin{itemize}[leftmargin=1.3em,itemsep=1pt]
\item \emph{central}: source $\colA$ and target seed $\colB$ (\Cref{def:columns});
\item \emph{external}: a plausible \emph{response} $\Resp(\colA)$ generated for
$\colA$ in its language, and $\Resp(\colB)$ for $\colB$ in its language, by the
demand-driven next-utterance of \Cref{thm:steer}.
\end{itemize}
The relaxation runs across all four: $\colA\leftrightarrow\colB$ (translation),
$\colA\!\to\!\Resp(\colA)$ and $\colB\!\to\!\Resp(\colB)$ (each unit generates the
conversation it would start), and $\Resp(\colA)\leftrightarrow\Resp(\colB)$ (do
the two start the same conversation?).
\end{construction}

\begin{theorem}[The four-column criterion]
\label{thm:four}
The four-column system is quiescent iff \textup{(a)} $\colA,\colB$ are quiescent
(same meaning) and \textup{(b)} $\Resp(\colA),\Resp(\colB)$ are quiescent (same
invited response). Quiescence of the whole certifies identity of meaning by the
\emph{route}---the responses agree---not merely by the endpoints---the content
agrees---and is computed without reading any meaning.
\end{theorem}

\begin{proof}
Quiescence of the system is the simultaneous vanishing of all four directions'
cross-demands (\Cref{def:relax}), which is exactly (a) and (b). Each is identity
of meaning of its pair by \Cref{thm:quiescence}, computed blind. (a) is identity
of the units; (b) is identity of what they propagate---the responses they
demand. The conjunction certifies that the two units \emph{start the same
conversation}, which by \Cref{thm:response,thm:steer} is the propagated meaning.
\end{proof}

\begin{theorem}[Route-audit: the four-column system catches the false friend]
\label{thm:false-friend}
Let $\colA,\colB$ be content-aligned units that nonetheless differ in meaning---a
\emph{false friend} (aligned endpoints, divergent interior). Then the central
relaxation may rest near quiescence on content while the external relaxation does
not: $\Resp(\colA)$ and $\Resp(\colB)$ have positive cross-demand, so the
four-column system is non-quiescent and the difference is exposed. An
endpoint-only criterion (content match) does not detect it; the response columns
do.
\end{theorem}

\begin{proof}
A false friend has $\colA,\colB$ landing on locally matched cells (small central
content gap) but different resting cuts (the meanings differ, by hypothesis). By
\Cref{thm:quiescence} the meaning-difference is a positive cross-demand; it
surfaces in what the units propagate, i.e.\ in the responses they invite, so
$\Resp(\colA)\neq\Resp(\colB)$ in meaning and
$\dem_{\Resp(\colA)\to\Resp(\colB)}>0$ (\Cref{thm:quiescence}). The four-column
demand is then positive (\Cref{thm:four}) though the central content gap is
small. The response layer audits the route through the propagated meaning, which
the endpoint content cannot.
\end{proof}

\begin{theorem}[Comprehension-free sufficiency]
\label{thm:compfree}
The four-column translator reaches the result of a meaning-dependent translator
without comprehension: it renders a unit by the response it invites (the
propagated meaning), exactly as a meaning-dependent translator does, while never
reading any meaning. The responses need not be correct---there is no correct
response (\Cref{thm:receiver})---only plausible propagations; by path-independence
which one is generated does not matter, only whether the response columns reach
quiescence.
\end{theorem}

\begin{proof}
A meaning-dependent translator renders ``How are you?'' as a greeting because it
routes through the response the utterance invites (\Cref{rem:criterion}). The
four-column system does the same by \Cref{thm:four}(b): quiescence of the
response columns is identity of the invited responses, which is identity of the
propagated meaning---the meaning-dependent translator's target. All quantities
are cross-demands, computed blind (\Cref{thm:quiescence}). The responses enter
only through their pairwise cross-demand; two plausible responses to the same
unit reverse-translate to the same demand on the other side, so the particular
response is immaterial---only the convergence of the response columns is used.
\end{proof}

\begin{remark}[Two columns for the statement, four for its propagation]
\label{rem:fourgen}
Two columns translate the \emph{statement} (its content relaxed against the
target); four translate the statement \emph{and what it propagates} (its meaning,
externalised as the response it invites). The construction generalises: a second
response turn is six columns, a tree of turns is the full conversation. The
minimal propagation-aware translator is the four-column one, and it is exactly
the upgrade from an endpoint-audit (does the content match?) to a route-audit
(does it start the same conversation?). No theorem depends on the reading.
\end{remark}

\begin{figure*}[t]
\centering
\includegraphics[width=\textwidth]{figures/panel_4.png}
\caption{\textbf{Route-audit, names, and the master equivalence.}
\textbf{(A)}~Central (content) demand versus response demand: proper pairs (blue)
sit at the origin, false friends (orange) have near-zero central demand---invisible
to an endpoint check---but large response demand, caught by the four-column
route-audit (\Cref{thm:false-friend}). \textbf{(B)}~A compound's cut $\sigma(U)$
(blue, flat at $2\flo$) versus the composition of its constituents' cuts (orange,
rising): names are non-compositional (\Cref{thm:non-derivable}).
\textbf{(C)}~The realised floor $\min_v\sigma$ versus $\flo$: positive for
$\flo>0$, vanishing (a sharp cut, red) only at $\flo=0$---the master equivalence
(\Cref{thm:master}). \textbf{(D)}~The same non-compositionality in three
coordinates $(W,\sigma(U),\text{composition})$.}
\label{fig:panel4}
\end{figure*}

% ============================================================================
\section{The master equivalence, scope, and conclusion}
\label{sec:master}
% ============================================================================

\begin{theorem}[Master equivalence]
\label{thm:master}
For a contact graph in a non-completable whole, the following are equivalent:
\begin{enumerate}[label=\textup{(\roman*)},leftmargin=2.0em,itemsep=1pt]
\item there is an identifiable item---a proper part told from the rest at
positive cost;
\item the floor is positive, $\flo>0$ (\Cref{thm:floor});
\item individuation is by negation without selector (\Cref{thm:negation}),
identity is the invariant minimum cut and never a point---the least sufficient
unknowability (\Cref{thm:region,thm:lsu}), meaning is unextractable
(\Cref{thm:unextractable}), processes are monotone non-returning residue chains
(\Cref{thm:propagator,thm:monotone}), uncertainty is distance to the floor
(\Cref{thm:replace}), and proper translation is the bidirectional relaxation to
quiescence, identity of meaning certified blind at positive residual of form
(\Cref{thm:quiescence,thm:form});
\item the perfect, traceless, point-meaning---a residue-free reading of a
thing---does not exist (\Cref{thm:unextractable}).
\end{enumerate}
Each is the same fact: telling a thing from the rest of a non-completable whole
costs a strictly positive, irreducible residue, which is at once the thing's
identity and what it propagates.
\end{theorem}

\begin{proof}
$(\mathrm{i})\Rightarrow(\mathrm{ii})$: an identifiable proper part has positive
separation cost (\Cref{thm:floor-derived,thm:floor}).
$(\mathrm{ii})\Rightarrow(\mathrm{iii})$: each listed structure is proved from
$\flo>0$ in the cited theorem.
$(\mathrm{iii})\Rightarrow(\mathrm{iv})$: \Cref{thm:unextractable} is among the
consequences. $(\mathrm{iv})\Rightarrow(\mathrm{i})$: if no residue-free reading
exists then every individuation costs $>0$, i.e.\ some thing is told from the
rest at positive cost---an identifiable item. The cycle closes.
\end{proof}

\subsection{Computational validation}

Because every object is a finite weighted graph and every minimum cut is
computable exactly by maximum flow \citep{fordfulkerson1956,menger1927}, each
theorem is checkable on explicitly constructed and randomly drawn instances. A
self-contained suite---an exact Edmonds--Karp max-flow / minimum-cut backend with
no external library, run under a fixed seed---implements one check per result and
records its outcomes. Across twenty categories it performed $4696$ individual
checks: the floor on every item of hundreds of random contact graphs; the
involution of negation; the invariance of the minimum cut under relabelling; the
region-valuedness of identity (the minimiser side non-singleton); the alignment
bound $\al\ge\flo/\Om$; the strict monotonicity of the committed record; the
residue--propagator equivalence; the entropy and information readings; the
equivalence of quiescence and an identical resting cut; the positive
form-residual; the relaxation dichotomy with decline on non-halting;
receiver-relativity; the four-column route-audit catching a false friend that an
endpoint check misses; the non-derivability and non-compositionality of names;
contextual interchangeability; the master floor / no-sharp-cut equivalence; and
the inconsistency of the self-name diagonal. \emph{All $4696$ checks passed}, with
minimum cuts computed exactly. The suite is a check against error and a display of
the quantitative shape of each result, not a substitute for the proofs, which hold
for every finite weighted graph in a non-completable whole.

\subsection{Scope}

Every result is a theorem about finite weighted graphs, provable from
\Cref{def:graph} and the derived floor \Cref{thm:floor}, with minimum cuts
computable exactly by maximum flow \citep{fordfulkerson1956,menger1927}; the
theorems are therefore checkable on explicitly constructed instances. The system
is qualitative-structural: it establishes \emph{that} individuation is a
positive-cost cut, \emph{that} identity is the least sufficient unknowability,
\emph{that} meaning is unextractable yet shareable, \emph{that} quiescence
constitutes identity of meaning, and \emph{that} a four-column relaxation
certifies it by route. It does not compute numerical cut weights, demands, or
alignments for any concrete language, which would require a concrete weighting.
One hypothesis beyond the graphs is used: that the whole is non-completable
(\Cref{def:noncompletable})---an order condition, not a cardinality posit, and
exactly what forces the floor; absent it, the floor may vanish and the
development is not compelled. Interpretive readings---of cells as meanings, of
residue as propagated significance, of the medium as the world, of receivers as
agents, of the response layer as a conversation---are confined to remarks on
which no theorem depends, and a reasonable reader will recognise that finitely
many bounded parts of many kinds satisfy the premise; we prove everything of the
abstract graph alone, and material that could not be proved at this standard has
been omitted.

\subsection{Conclusion}

From a single derived fact---that telling a part from the rest of a
non-completable whole costs a strictly positive, irreducible floor---an entire
theory of meaning follows. A thing's identity is what it is \emph{not},
committed at least sufficient cost: a region, never the thing itself, for to be
the thing itself one would have to exhaust the inexhaustible. A \emph{name} is
that cut handed back as a token---an irreducible, content-free stash of what a
thing is not---and the namer, with the largest complement, names itself last and
only approximately, so language begins at the periphery and converges on the
source it never reaches. Meaning is that
identity together with what it propagates as a cause; it cannot be read, because
reading it would be being it, but identity of meaning can be \emph{detected}, by
the absence of any difference that would propagate. Translation is then not the
transfer of a content but the relaxation of two sides to semantic
quiescence---the silence of all demands---reached at positive residual of form,
and where no fixed point exists, declined rather than faked. Conversation is the
same propagation kept open: a response is forced precisely because no meaning can
be extracted, and is steered by the demand, not by comprehension. The
four-column translator externalises the propagation as the response a unit
invites, turning translation into the test of whether two utterances start the
same conversation---a route-audit a meaning-blind mechanism passes and a false
friend fails. Uncertainty, throughout, is not the obstacle to meaning but its
condition: the floor that forbids perfect alignment is the floor that makes there
be anything to say, and a finite system that does not need to understand is
sufficient, because identity of meaning is the silence of all demands, and that
silence can be reached, and recognised, blind.

% ============================================================================
\bibliographystyle{plainnat}
\bibliography{references}
% ============================================================================

\end{document}